% Generated from tex/chapters/ch04/00_chapter_opening.tex for the SWEBOK structure tree
\section{第04章 ソフトウェア構築}
\noindent\textbf{この資料の主張}

ソフトウェア構築は、単にプログラムを書く作業ではない。要求と設計を、実行でき、検証でき、保守できるソフトウェアへ変える中心的な活動である。構築には、コーディング、検証、単体テスト、統合テスト、デバッグ、再利用、統合、依存関係管理、性能調整、開発環境の選択が含まれる。よい構築は、複雑さを抑え、変化を受け入れ、欠陥を早く見つけ、既存資産を適切に使い、標準に沿って進めることで実現する。

\begin{plainbox}
\textbf{この資料の読み方}

専門用語は原則として日本語で示し、必要に応じて英語を括弧内に併記する。英語名は、元資料やツール上の名称を確認したいときの手がかりとして残す。初学者が前提知識なしに読めるように、各節では「何のための概念か」「実務では何を見ればよいか」を重視する。
\end{plainbox}
